\documentclass{article}

\usepackage{amsmath,amssymb}
\usepackage{xurl}
\usepackage[hidelinks]{hyperref}

% Shared commands for notes in docs/paper-gaps/.

\DeclareUnicodeCharacter{2080}{\ifmmode{}_{0}\else\textsubscript{0}\fi}
\DeclareUnicodeCharacter{2081}{\ifmmode{}_{1}\else\textsubscript{1}\fi}
\DeclareUnicodeCharacter{2082}{\ifmmode{}_{2}\else\textsubscript{2}\fi}
\DeclareUnicodeCharacter{2099}{\ifmmode{}_{n}\else\textsubscript{n}\fi}

\newcommand{\Lean}{\textsc{Lean}}
\ExplSyntaxOn
\NewDocumentCommand{\leanid}{m}{
  \ifmmode
    \textsf{\detokenize{#1}}
  \else
    \group_begin:
    \urlstyle{sf}
    \tl_set:Nn \l_tmpa_tl {#1}
    \tl_replace_all:Nnn \l_tmpa_tl {\_} {_}
    \exp_args:NV \path \l_tmpa_tl
    \group_end:
  \fi
}
\ExplSyntaxOff
\providecommand{\tr}{\operatorname{tr}}
\newcommand{\tnleanIssueBase}{https://github.com/LionSR/TNLean/issues/}
\newcommand{\tnleanPrBase}{https://github.com/LionSR/TNLean/pull/}
\newcommand{\tnleanDocBase}{https://sirui-lu.com/TNLean/docs/}
\newcommand{\ghissue}[1]{\href{\tnleanIssueBase#1}{GitHub issue \##1}}
\newcommand{\ghpr}[1]{\href{\tnleanPrBase#1}{PR \##1}}
\newcommand{\doclink}[1]{\href{\tnleanDocBase#1.html}{\path{#1}}}

% Dirac notation and the identity operator, as in blueprint/src/macros/common.tex.
% \providecommand keeps note-local definitions authoritative.
\usepackage{dsfont}
\providecommand{\ket}[1]{|#1\rangle}
\providecommand{\bra}[1]{\langle #1|}
\providecommand{\braket}[2]{\langle #1 | #2 \rangle}
\providecommand{\ketbra}[2]{|#1\rangle\!\langle #2|}
\providecommand{\Id}{\mathds{1}}
\providecommand{\one}{\mathds{1}}
\providecommand{\C}{\mathbb{C}}
\providecommand{\MN}[1]{M_{#1}(\C)}
\providecommand{\MD}{M_D(\C)}

% External referencing for paper sources.  The build helper writes sanitized
% auxiliary files under build/paper-aux/ and excludes duplicated source labels.
\usepackage{xr}

% Import a generated paper auxiliary file with a caller-supplied label prefix.
% The candidate paths support builds from docs/paper-gaps/, the repository root,
% and build/paper-aux/ itself.
\newcommand{\importpaperaux}[2]{%
  \IfFileExists{../../build/paper-aux/#2.aux}{%
    \externaldocument[#1]{../../build/paper-aux/#2}%
  }{%
    \IfFileExists{build/paper-aux/#2.aux}{%
      \externaldocument[#1]{build/paper-aux/#2}%
    }{%
      \IfFileExists{#2.aux}{%
        \externaldocument[#1]{#2}%
      }{}%
    }%
  }%
}

% General external reference macros
\newcommand{\paperref}[2]{\ref{#1-#2}}
\newcommand{\papereqref}[2]{(\ref{#1-#2})}

% CPSV16 source (1606.00608 / MPDO-22-12-17-2.tex) external document and shortcuts
\importpaperaux{cpsv16-}{1606.00608/MPDO-22-12-17-2-tnlean}
\newcommand{\cpsvref}[1]{\paperref{cpsv16}{#1}}
\newcommand{\cpsveqref}[1]{\papereqref{cpsv16}{#1}}

% Verdict marker: \gapnote{<kind>}{<status>}. Kind is the policy
% classification (clarification, local-correction, scope-restriction,
% unfaithful, false-source, open-gap); status is open, wip (elimination
% underway), resolved, or historical. Machine-read by the site index;
% prints a verdict line.
\newcommand{\gapnote}[2]{\par\noindent\textbf{Verdict:} #1 (#2).\par}


\title{The Endpoint Rank-Product Exponent in MPU Index Well-Definedness}
\author{}
\date{2026-08-12}

\begin{document}
\maketitle
\gapnote{local-correction}{resolved}

\section{The printed exponent}

Let $d$ be the one-site physical dimension, and let $r_k$ and $\ell_k$ denote
the right and left source-cut ranks of the $k$-site blocked tensor. In the
proof of Proposition~IV.2 of Ref.~\cite{Cirac2017MPU}, the rank product is
printed as
\begin{align}
    r_k\ell_k
        &=d^k.
    \label{eq:mpu_blocking_rank_product_exponent_printed_identity}
\end{align}
The blocked physical dimension is
\begin{align}
    q
        &=d^k.
    \label{eq:mpu_blocking_rank_product_exponent_blocked_dimension}
\end{align}

\section{Consistent rank-product identity}

Theorem~III.8 of Ref.~\cite{Cirac2017MPU} states that a simple tensor of
physical dimension $q$ has source-cut ranks satisfying
\begin{align}
    r\ell
        &=q^2.
    \label{eq:mpu_blocking_rank_product_exponent_simple_tensor_identity}
\end{align}
Applying
(\ref{eq:mpu_blocking_rank_product_exponent_simple_tensor_identity}) to the
$k$-site block and using
(\ref{eq:mpu_blocking_rank_product_exponent_blocked_dimension}) gives
\begin{align}
    r_k\ell_k
        &=q^2
    \label{eq:mpu_blocking_rank_product_exponent_blocked_rank_product}\\
        &=(d^k)^2
    \label{eq:mpu_blocking_rank_product_exponent_substituted_dimension}\\
        &=d^{2k}.
    \label{eq:mpu_blocking_rank_product_exponent_corrected_identity}
\end{align}
The same square appears in the index discussion and in the later blocking
argument of Ref.~\cite{Cirac2017MPU}. Thus
(\ref{eq:mpu_blocking_rank_product_exponent_printed_identity}) contains a
typographical exponent error. The consistent endpoint identity is
(\ref{eq:mpu_blocking_rank_product_exponent_corrected_identity}).

\section{Formal treatment}

The simultaneous saturation theorem
\leanid{MPOTensor.blockingRanks_eq_pow_mul_of_products} does not derive the
rank-product identity. It assumes the endpoint identities
\begin{align}
    r_{k_0}\ell_{k_0}
        &=d^{2k_0},
    \label{eq:mpu_blocking_rank_product_exponent_initial_endpoint_identity}\\
    r_k\ell_k
        &=d^{2k}.
    \label{eq:mpu_blocking_rank_product_exponent_final_endpoint_identity}
\end{align}
It combines these identities with independently proved blocking upper bounds.
The focused accessors are
\leanid{MPOTensor.rightRank_blockTensor_eq_pow_mul_of_products} and
\leanid{MPOTensor.leftRank_blockTensor_eq_pow_mul_of_products}. A later
application may derive the endpoint identities from the formal counterpart
of Theorem~III.8 of Ref.~\cite{Cirac2017MPU}, under that theorem's own
hypotheses.

\section{Verdict}

The exponent in the proof of Proposition~IV.2 of
Ref.~\cite{Cirac2017MPU} is corrected from $d^k$ to $d^{2k}$. This resolved
local correction adds no hypothesis beyond the rank-product identity used by
the proof. The formal saturation theorem nevertheless retains that identity
as an explicit hypothesis rather than deriving it from blocking alone.

\bibliographystyle{alpha}
\bibliography{references}

\end{document}
